\documentclass[11pt,a4paper]{article}

% ---------- Packages ----------
\usepackage[utf8]{inputenc}
\usepackage[T1]{fontenc}
\usepackage[a4paper,margin=2.2cm]{geometry}
\usepackage{xcolor}
\usepackage[hidelinks]{hyperref}
\usepackage{parskip}

% ---------- Style (matches resume) ----------
\definecolor{accent}{RGB}{38,86,166}
\hypersetup{colorlinks=true,urlcolor=accent}
\pagestyle{empty}
\setlength{\parskip}{8pt}

\begin{document}

% ---------- Header ----------
{\LARGE\bfseries\color{accent} Mehmet Çağrı Ekici}\par
\vspace{2pt}
{\small
\href{mailto:mehmetcagriekici@gmail.com}{mehmetcagriekici@gmail.com} \,\textbullet\,
+90 536 477 17 89 \,\textbullet\,
Ankara, Turkey \,\textbullet\,
\href{https://github.com/mehmetcagriekici}{github.com/mehmetcagriekici} \,\textbullet\,
\href{https://linkedin.com/in/mehmetcagriekici}{linkedin.com/in/mehmetcagriekici}
}

\vspace{10pt}
\hrule
\vspace{14pt}

Dear Olive Tree Team,

I'm applying for the Junior React Developer position. I'll be upfront about what my profile is: I am a backend-first developer with genuine hands-on React experience from side projects. I want to explain why I think that combination fits this role better than it might sound at first.

The backend-first part: since July 2025 I have been focused on backend engineering. I completed Boot.dev's Backend Developer Path twice --- once in Python and Go, again in TypeScript --- and I build HTTP servers and clients myself. Your posting says the developer will connect frontend components to APIs to display real-time updates. That is where this background pays off: when I wire a component to an endpoint, I understand both sides of the fetch call --- what the server does with the request, why the response is shaped the way it is, and what to do when something breaks. For dashboards whose whole job is showing data clearly, knowing where the data comes from is not a side skill.

The React part is real, not resume decoration. I built \href{https://github.com/mehmetcagriekici/rotjournal}{rotjournal}, a journaling app in React 18 and TypeScript with Redux Toolkit for state, a rich-text editor, and react-hook-form. I built several Next.js applications with React 19 --- among them a daily Bible study platform using Redux Toolkit, React Query, react-hook-form, Material UI, and Supabase for auth and data. All of them use Vite or Next tooling, npm, and Git. Before frameworks, I completed many Frontend Mentor challenges in plain HTML, CSS, and JavaScript, rebuilding layouts from design files --- work where spacing and visual consistency are the entire job, which is the attention to detail your posting asks for. These are side projects, so my React would need a short warm-up after months of backend work --- days, not months --- but the component model, state management, and forms are things I have used, not just read about.

There is also a reason this role in particular caught my eye. Your firm's principle that ``memory is not a management system'' describes my own portfolio better than I could have. rotjournal is a tool for writing things down instead of trusting memory. My current backend project, BlightSanest, is a platform where users keep records across any domain and query them later, so that nothing important depends on who remembers it. Building dashboards that log decisions and ownership is the same idea, applied to teams --- work I would actually care about, not just accept.

Two more honest notes. This would be my first professional role, and code review with an experienced team is exactly what I am applying for. And on remote self-discipline, I can offer evidence instead of promises: I have studied and built nearly every day since July 2025, verifiable at \href{https://www.boot.dev/u/callsower}{boot.dev/u/callsower}.

I would love to show your team how I work. Thank you for your time.

Sincerely,

\textbf{Mehmet Çağrı Ekici}

\end{document}
